\section{Onde aprofundar}

\subsection*{Livros}
\small
\begin{longtable}{@{}p{1.8cm} p{4.4cm} p{2.6cm} p{4cm}@{}}
\toprule
Tema & Livro & Autor(es), ano & Por que ler \\
\midrule \endhead
Desempenho & \citetitle{livro:oaks2020}, 2ª ed. & \citeauthor{livro:oaks2020},
  \citeyear{livro:oaks2020} & JIT, GC, memória, threads e metodologia de
  teste de desempenho, com foco em Java~8 e 11 \\
Desempenho & \citetitle{livro:evans2018} & \citeauthor{livro:evans2018},
  \citeyear{livro:evans2018} & Abordagem quantitativa: logs de GC, JIT,
  ferramentas e os perigos dos microbenchmarks \\
Concorrência & \citetitle{livro:goetz2006} & \citeauthor{livro:goetz2006},
  \citeyear{livro:goetz2006} & Visibilidade, publicação segura e
  utilitários de \texttt{java.util.concurrent} \\
\bottomrule
\end{longtable}
\normalsize

Os dois livros de desempenho são anteriores ao ZGC geracional e ao cache
AOT; use-os pelos conceitos e confira opções e padrões nos JEPs e na
documentação do JDK~25.

\subsection*{Artigos}
Da série \emph{JVM Anatomy Quarks}, de Aleksey Shipilëv:
\citetitle{artigo:shipilev-uncommon-traps},
\citetitle{artigo:shipilev-megamorphic} e
\citetitle{artigo:shipilev-safepoint}. Sobre o modelo de memória:
\citetitle{artigo:shipilev-jmm}.

\subsection*{Documentação}
\citetitle{doc:jvms25} (cap.~5), \citetitle{doc:jls25} (cap.~17),
\citetitle{doc:hotspot-perf}, \citetitle{doc:gc-collectors} e
\citetitle{doc:java-cmd}.
